\section{From the regular point to the original representation}
\label{sec:conclusion}

We have a fixed-determinant trace domain $D$ containing both the
normalized original point and the regular point $x$ of
Proposition~\ref{prop:regular-point}.  The latter point is now on
Hecke support over $\Q$.  To put the whole domain on that support,
we compare both quotients of the global deformation ring in which
the determinant is allowed to vary.  The decisive fact is that this
larger ring is regular at~$x$.

\subsection{The full adjoint tangent space}

\begin{proposition}\label{prop:tangent-bound}
Let $r_x:G_{\Q,S}\to\GL_2(E_x)$ be odd, absolutely irreducible,
and classical cuspidal up to Tate twist.  Suppose its restriction
to $G_{\Q_2}$ is absolutely irreducible and de Rham with distinct
Hodge--Tate weights, and its Weil--Deligne representation is trivial
on wild inertia.  Here $S$ is any finite set containing $2$ and all
ramified primes.  Then
\[
 \dim_{E_x}H^1(G_{\Q,S},\ad r_x)\leq3.
\]
Consequently, if $R_{\mathrm{glob}}$ is the universal global
determinant deformation ring with varying determinant for the
semisimplified residual data of $r_x$, and $\mathfrak x$ is its
coefficient prime, then
\[
 \edim (R_{\mathrm{glob}})_{\mathfrak x}\leq3.
\]
Both assertions permit scalar or reducible semisimplified residual
data and finite extension of the coefficient ring.
\end{proposition}

\begin{proof}
Set $W=\ad r_x=\End_{E_x}(r_x)$, the full four-dimensional
adjoint representation. Tameness of the Weil--Deligne parameter will
give the CM qualification for adjoint Selmer vanishing; local
irreducibility and distinct de Rham weights will give the local
cohomology bound. For a finite place $v$, the finite local
condition is
\[
 H^1_f(\Q_v,W)=
 \begin{cases}
  \ker\bigl(H^1(\Q_v,W)\to H^1(I_v,W)\bigr),&v\ne2,\\
  \ker\bigl(H^1(\Q_2,W)\to
        H^1(\Q_2,W\otimes_{\Q_2}B_{\mathrm{cris}})\bigr),&v=2.
\end{cases}
\]
Here $B_{\mathrm{cris}}$ is Fontaine's crystalline period ring.
The global group $H^1_f(\Q,W)$ imposes these conditions at all
finite places.  Positive-degree cohomology at the real place is
zero over $E_x$.

Newton--Thorne's theorem \cite[Theorem~5.4]{NewtonThorne2023},
specialized to $\GL_2/\Q$, states that for a regular algebraic
cuspidal automorphic representation $\pi$ and any prime $p$,
\[
 H^1_f(\Q,\ad r_{\pi,p})=0
\]
provided either $\pi$ is non-CM, or its CM field $K$ is not
contained in $\Q(\zeta_{p^\infty})$.  This statement uses the
full adjoint and has no residual irreducibility hypothesis.
Our classical representation is regular algebraic, including
when its classical weight is two, and twists do not change its
adjoint.  It remains only to verify the CM qualification.

If its CM field were contained in $\Q(\zeta_{2^\infty})$, the
corresponding quadratic self-twist would restrict to a nontrivial
wildly ramified character at $2$.  Indeed the three quadratic
subfields are $\Q(i)$, $\Q(\sqrt2)$, and $\Q(\sqrt{-2})$;
each ramifies at $2$, and a quadratic ramified character there
is necessarily wild.  The Weil--Deligne parameter would then
be isomorphic to its twist by that character.  On a wild inertia
element where the character is $-1$, these two parameters act
as $I$ and $-I$, respectively, by the hypothesis on~$x$.
This is impossible.  Newton--Thorne therefore gives
\begin{equation}\label{eq:global-adjoint-selmer}
 H^1_f(\Q,W)=0.
\end{equation}

We next compute the local quotients.  The trace pairing identifies
$W^*$ with $W$, so local Tate duality identifies the dual of
$H^2(\Q_v,W)$ with $\Hom_{G_{\Q_v}}(r_x,r_x(1))$.
At $2$, a nonzero such map is an isomorphism by absolute
irreducibility; its determinants would force $\eps^2=1$ on
$G_{\Q_2}$. Thus $H^2(\Q_2,W)=0$, including in weight two.
At $v\ne2$, local--global compatibility
\cite[Theorem~A]{Carayol1986} and genericity of each local factor
of a cuspidal automorphic representation give the same vanishing:
the genericity criterion
\cite[Lemma~1.1.3]{Allen2016} says that its Frobenius-semisimple
Weil--Deligne parameter $U$ satisfies
$\Hom_{\mathrm{WD}}(U,U(1))=0$.
A Galois morphism would induce such a morphism, also after
Frobenius semisimplification.  Tate twists do not affect these
adjoint computations.  This is the non-dyadic argument of
\cite[Lemma~5.2]{Dim}.

Write $h^j_v=\dim_{E_x}H^j(\Q_v,W)$.
For $v\ne2$, local Euler characteristic gives $h^1_v=h^0_v$,
whereas
\[
 H^1_f(\Q_v,W)=W^{I_v}/(\operatorname{Frob}_v-1)W^{I_v}
\]
also has dimension $h^0_v$.  Hence the local quotient is zero.
At $2$, Schur's lemma gives $h^0_2=1$, and the Euler
characteristic formula gives $h^1_2=1+\dim W=5$.
The filtration degrees of the filtered de Rham module
$D_{\mathrm{dR}}(W)$ are
$0,0,h,-h$ for some nonzero integer $h$.  The Bloch--Kato
dimension formula \cite[Corollary~3.8.4]{BlochKato1990} gives
\[
 \dim H^1_f(\Q_2,W)
 =h^0_2+
   \dim D_{\mathrm{dR}}(W)/\operatorname{Fil}^0=1+1=2.
\]
We have proved
\begin{equation}\label{eq:local-adjoint-quotients}
 \dim H^1(\Q_v,W)/H^1_f(\Q_v,W)
 =\begin{cases}0,&v\ne2,\\3,&v=2.\end{cases}
\end{equation}
The kernel of the restriction map
\[
 H^1(G_{\Q,S},W)\longrightarrow
 \bigoplus_{v\in S} H^1(\Q_v,W)/H^1_f(\Q_v,W)
\]
lies in $H^1_f(\Q,W)$, because classes in its source are already
unramified outside $S$.  Equations~\eqref{eq:global-adjoint-selmer}
and~\eqref{eq:local-adjoint-quotients} now give the bound three.

For the embedding-dimension assertion, apply the characteristic-zero
tangent comparison of \cite[Proposition~4.1]{Dim} to
$R_{\mathrm{glob}}$. Its trace and determinant values topologically
generate the ring, and the specialization at $\mathfrak x$ is
absolutely irreducible. The comparison bounds the embedding dimension by
$\dim_{E_x}H^1(G_{\Q,S},W)$, with no condition on the
semisimplified residue. The continuity needed here follows because
$R_{\mathrm{glob}}/\mathfrak x^2$ is finite over $\OO$:
$R_{\mathrm{glob}}/\mathfrak x$ is finite over $\OO$, and
$\mathfrak x/\mathfrak x^2$ is finite over that quotient.
Finite extension of $\OO$ adds no tangent directions, since
characteristic-zero derivations kill its fraction field by
separability over $\Q_2$. The preceding cohomological bound therefore
gives the asserted embedding dimension.
\end{proof}

\subsection{A regular point determines the supporting component}

The following elementary lemma isolates the use of the
four-dimensional Hecke theorem.  It explains why the tangent
calculation is needed in addition to the existence of a modular point.

\begin{lemma}\label{lem:regular-support}
Let $R_0$ be a complete Noetherian local $\OO$-algebra, and let
$A_0$ and $D_0$ be quotients of $R_0$, with $D_0$ a domain.
Suppose every irreducible component of $\Spec A_0$ has dimension
four.  Let $x_0:R_0\to\OO'$ be a continuous coefficient point
factoring through both quotients, where $\OO'/\OO$ is a finite
extension of integer coefficient rings.  If
\[
 \edim(R_0)_{\mathfrak x_0}\leq3,
 \qquad \mathfrak x_0=\ker x_0,
\]
then $R_0\to D_0$ factors through $A_0$.
\end{lemma}

\begin{proof}
Choose a minimal prime $P$ of $A_0$ contained in its prime
$\mathfrak x_0$.  The complete local domain $A_0/P$ is
catenary.  Moreover $A_0/\mathfrak x_0$ is an $\OO$-submodule
of $\OO'$, hence finite over $\OO$, and it contains $\OO$.
It therefore has dimension one.  The dimension formula gives
\[
 \dim(A_0/P)_{\mathfrak x_0}
 =\dim A_0/P-\dim A_0/\mathfrak x_0=4-1=3.
\]
As $R_0$ surjects onto $A_0$, we obtain
\[
 3\leq\dim (A_0)_{\mathfrak x_0}
 \leq\dim (R_0)_{\mathfrak x_0}
 \leq\edim (R_0)_{\mathfrak x_0}\leq3.
\]
Thus $(R_0)_{\mathfrak x_0}$ is a regular local ring of dimension
three, in particular a domain.  Its quotient
$(A_0)_{\mathfrak x_0}$ has the same dimension, so the localized
kernel $\ker(R_0\to A_0)_{\mathfrak x_0}$ is zero: a nonzero
ideal in this regular local domain has quotient of smaller
dimension.

If $a\in\ker(R_0\to A_0)$, there is consequently
$s\notin\mathfrak x_0$ with $sa=0$ in $R_0$.
The kernel of $R_0\to D_0$ is contained in $\mathfrak x_0$,
so the image of $s$ in $D_0$ is nonzero.  Since $D_0$ is a
domain, the image of $a$ must vanish.  This proves the factorization.
\end{proof}

\begin{proof}[Proof of Theorem~\ref{thm:main}]
Apply Proposition~\ref{prop:family} to the original representation,
after its indicated continuous cyclotomic normalization and finite
coefficient extension.  It supplies the horizontal trace domain
$D$ with the normalized original coefficient point.
Propositions~\ref{prop:diagonal}, \ref{prop:positive-rank},
and~\ref{prop:regular-point} supply a coefficient point $x$ of
$D$ satisfying all the hypotheses of
Proposition~\ref{prop:tangent-bound}.

Choose an odd $N\geq5$ sufficiently divisible that it contains
every odd ramified prime of the original family and that the
classical representation underlying $r_x$ has level dividing
$2^aN$ for some $a$.  Increase the allowed global ramification
set to the primes dividing $2N$; the existing family remains
unramified at any newly added primes.  By
Propositions~\ref{prop:hecke-level} and~\ref{prop:hecke-twists},
$r_x$ gives a continuous point of $\mathbb T_2(N)$ with the
required Frobenius values.

Fix the coefficient ring $\OO$ of $D$ and the residual embedding
selected by~$x$.  Let $A$ be the corresponding residue factor
of $\OO\otimes_{\Z_2}\mathbb T_2(N)$, and let
$R_{\mathrm{glob}}$ be the universal global determinant deformation
ring for this residual determinant with \emph{varying} determinant.
Let $\OO_x$ be a coefficient ring for $x$. The determinant on $D$
and the Hecke determinant give a commutative diagram whose top
arrows are continuous surjections:
\[
 \begin{array}{ccc}
   &R_{\mathrm{glob}}&\\[-2pt]
   \swarrow&&\searrow\\[-2pt]
   D&&A\\[-2pt]
   \searrow&&\swarrow\\[-2pt]
   &\OO_x.&
 \end{array}
\]
The lower arrows are the coefficient point $x$; equality follows
from its Frobenius identities and Chebotarev.  The quotient $D$
has fixed determinant $\chi$, whereas $R_{\mathrm{glob}}$ and
$A$ retain determinant variation.  Surjectivity onto $A$ follows
from topological generation by the Hecke elements and compactness,
as in \cite[Proposition~2.3]{Dim}; surjectivity onto $D$ is its
definition as a trace-image quotient.

Every component of $A$ has dimension four by
\cite[Theorem~1.1]{Dim} and finite coefficient extension.
Proposition~\ref{prop:tangent-bound} gives
$\edim(R_{\mathrm{glob}})_{\mathfrak x}\leq3$.
Lemma~\ref{lem:regular-support} therefore makes
$R_{\mathrm{glob}}\to D$ factor through $A$.
The normalized original coefficient point of $D$ is consequently
a point of the required Hecke algebra over~$\Q$.

Undo the normalization with Proposition~\ref{prop:hecke-twists}.
We obtain, over a finite extension $E'/E$, a continuous map
\[
 \lambda:\mathbb T_2(N)\longrightarrow\OO_{E'}
\]
whose values on $T_\ell$ and $\ell S_\ell$ are respectively
$\tr r(\operatorname{Frob}_\ell)$ and
$\det r(\operatorname{Frob}_\ell)$ for $\ell\nmid2N$.
The compact image of $r$ preserves an $\OO_E$-lattice, so all
these values lie in $\OO_E$.  Since the Hecke generators
topologically generate $\mathbb T_2(N)$ and $\OO_E$ is closed
in $\OO_{E'}$, the entire image of $\lambda$ lies in $\OO_E$.
The induced map to $\OO_E$ is continuous for its $2$-adic topology.
Our choice of $N$ includes every odd ramified prime of the
original $r$, which completes the proof.
\end{proof}
